\documentclass[12pt]{article}
%\include{amsfonts}
\usepackage{amssymb,amsmath,amsthm,latexsym,epsfig,euscript,multicol}
\usepackage{enumitem}
\usepackage[utf8x]{inputenc}
\usepackage{booktabs}
\usepackage{babel}[spanish]
\usepackage{multicol}
% \usepackage{enumerate}

% Caracteres especiales
\def\A{\mathbb{A}}
\def\C{\mathbb{C}}
\def \N{\mathbb{N}}
\def \P{\mathbb{P}}
\def \Q{\mathbb{Q}}
\def \R{\mathbb{R}}
\def \Z{\mathbb{Z}}
\def \sen{\textrm{sen}}

\theoremstyle{definition}
\newtheorem{ejer}{Ejercicio}
\newcommand{\bej}{\begin{ejer}}
\newcommand{\fej}{\end{ejer}}

\def\dt{\Delta t}
\def\dx{\Delta x}

\topmargin-2cm \vsize 29.5cm \hsize 21cm
\setlength{\textwidth}{16.75cm}\setlength{\textheight}{23.5cm}
\setlength{\oddsidemargin}{0.0cm}
\setlength{\evensidemargin}{0.0cm}


\begin{document}

%----------------------------------------------------------
%Encabezado:
\begin{enumerate}
    \item Los siguientes diccionarios corresponden soluciones básicas factibles de modelos de PL con objetivo de
    \textbf{maximizar}.
    Realizar una iteración de SIMPLEX. ¿Qué se observa?
        \begin{enumerate}[label=\alph*)]
        \item  \[
        \begin{array}{ccccccccc}
            x_2 & = & 5 & + & 2x_3 & - & x_4  & - & 3x_1 \\
            x_5 & = & 7 &   &      & - & 3x_4 & - & 4x_1 \\ \hline
            z   & = & 5 & + & x_3  & - & x_4  & - & x_1  \\
        \end{array}
        \]
        \item
        \[
        \begin{array}{ccccccccc}
            x_2 & = & 1 & - & 2x_1 & - & x_4   \\
            x_3 & = &     & - & x_1 & - & x_4\\ \hline
            z   & = & 4   & + & 2x_1 & - & x_4\\
        \end{array}
        \]
        \end{enumerate}
    \item El siguiente diccionario corresponde a una solución óptima de un modelo de PL con objetivo de
    \textbf{maximizar}.
    Si aumentamos el valor de $x_2$, ¿qué sucede con el valor de la f.o.? ¿Cuánto puede aumentar $x_2$ y seguir
    garantizando la no negatividad de las demás variables?
        \[
        \begin{array}{ccccccccc}
            w_1 & = & 3 & + & x_2  & - & 2w_2 & + & 7x_3 \\
            x_1 & = & 1 & - & 5x_2 & + & 6w_2 & - & 8x_3 \\
            w_3 & = & 4 & + & 9x_2 & + & 2w_2 & - & x_3  \\ \hline
            z   & = & 8 &   &      & - & w_2  & - & x_3  \\
        \end{array}
        \]
\end{enumerate}
\end{document}
